\begin{abstract}
Flood segmentation models for Sentinel-1 radar must work on new floods, yet they are usually tested on
images from the flood events they were trained on. We show that this hides the value of terrain. A U-Net
trained on the 11 events of Sen1Floods11 gains almost nothing from height above nearest drainage on that
dataset's own test split (+0.007 intersection over union) but +0.128 on its held-out Bolivia event. On 79
external Copernicus emergency activations, the same input raises mean event intersection over union from
0.574 to 0.635 for water and from 0.317 to 0.349 for flooded area, with water improved on 69 of the 79 events, and the
gain holds without a threshold and under the benchmark's pre- and post-event protocol. It needs the right
terrain for the place: elevation and slope recover only a quarter of it, its event-level gain vanishes when
the drainage input is shifted by 500 m, blurred, or taken from another flood, and a model trained on such mismatched
input gains nothing. A fine-tuned 101-million-parameter foundation model benefits about as much and,
without the input, does not match the 7.8-million-parameter U-Net that has it. Used as a physics-based loss
or a mask, the same terrain does nothing. Trained on 159 events, the flood-extent gain is no longer
detectable, although a pre-registered test did not confirm a steady decline with the number of events. With
few labelled floods, terrain is a cheap gain that same-event evaluations miss.
\end{abstract}

\begin{IEEEkeywords}
Flood mapping, synthetic aperture radar, terrain, domain generalization.
\end{IEEEkeywords}
